% ADD THIS HEADER TO ALL NEW CHAPTER FILES FOR SUBFILES SUPPORT

% Allow independent compilation of this section for efficiency
\documentclass[../CLthesis.tex]{subfiles}

% Add the graphics path for subfiles support
\graphicspath{{\subfix{../images/}}}

% END OF SUBFILES HEADER

%%%%%%%%%%%%%%%%%%%%%%%%%%%%%%%%%%%%%%%%%%%%%%%%%%%%%%%%%%%%%%%%
% START OF DOCUMENT: Every chapter can be compiled separately
\begin{document}

\chapter{RQ3: What Moral Motivations Drive Anti-Religious Hate Speech, and How Do They Differ Across Languages and Topics?}%
\label{chap:why}

Building on the moral foundation scoring framework introduced in Chapter~\ref{chap:methods}, this chapter investigates \textit{what moral motivations} underlie anti-religious hate speech — and whether those motivations systematically differ by language or by topic. 

Before interpreting individual patterns, we first verify that these bias scores carry meaningful signal. A one-way ANOVA across all three language groups (English, Chinese, Japanese) was conducted for each of the seven axes. As shown in Table~\ref{tab:rq3_anova}, all seven axes yield extremely significant differences ($p < 4.9 \times 10^{-324}$ for the five Moral Foundation axes, and $p < 10^{-200}$ for the two threat axes), with large effect sizes ($\eta^2$ ranging from 0.24 to 0.95). This confirms that the moral bias scores capture real, language-dependent variation that is worth investigating.

\begin{table}[htbp]
\centering
\caption{One-way ANOVA results for moral bias scores across English, Chinese, and Japanese corpora. All pairwise Tukey post-hoc comparisons are also significant ($p < 10^{-10}$) for all axes.}
\label{tab:rq3_anova}
\begin{tabularx}{\textwidth}{l X r r l}
\toprule
\textbf{Axis} & \textbf{Description} & \textbf{F-stat} & \boldmath$\eta^2$ & \textbf{\textit{p}-value} \\
\midrule
Harm             & Care(+) vs.\ Harm(−)                         & 32{,}068 & 0.949 & $< 4.9 \times 10^{-324}$ \\
Fairness         & Fairness(+) vs.\ Cheating(−)                 & 35{,}568 & 0.954 & $< 4.9 \times 10^{-324}$ \\
Loyalty          & Loyalty(+) vs.\ Betrayal(−)                  & 10{,}808 & 0.863 & $< 4.9 \times 10^{-324}$ \\
Authority        & Authority(+) vs.\ Subversion(−)              & 17{,}686 & 0.912 & $< 4.9 \times 10^{-324}$ \\
Sanctity         & Sanctity(+) vs.\ Degradation(−)              & 25{,}366 & 0.937 & $< 4.9 \times 10^{-324}$ \\
Realistic Threat & Safety(+) vs.\ Realistic Threat(−)           &    840   & 0.329 & $1.03 \times 10^{-297}$  \\
Symbolic Threat  & Cultural Cohesion(+) vs.\ Symbolic Threat(−) &    544   & 0.241 & $6.71 \times 10^{-206}$  \\
\bottomrule
\end{tabularx}
\end{table}

\paragraph{How to read bias scores.}
Before interpreting the figures, one semantic caveat deserves emphasis. As defined in Section~\ref{sec:rq3}, the bias score is the cosine alignment between a document and an axis vector anchored by two opposing vocabulary clusters in a shared multilingual embedding space. It therefore reflects \emph{which moral vocabulary the document activates}, not whether the speaker is praising or attacking the axis concept. A positive score on Sanctity, for instance, does not mean that the text treats religion as sacred — it means the document's language is lexically closer to purity/tradition vocabulary than to degradation vocabulary, which is fully consistent with a speaker invoking those virtue words in order to attack an outgroup (``we must protect our sacred tradition from this foreign cult'' would score positively on Sanctity precisely while being hostile).

Concretely, anti-religious discourse follows two distinct moral-argumentative grammars, both of which are hostile:
\begin{itemize}
    \item \textbf{Direct-condemnation grammar} — the speaker loads vice vocabulary (\textit{harm}, \textit{corruption}, \textit{cheating}, \textit{degradation}) directly onto the target, yielding a \emph{negative} bias on the relevant MFT axis. The target is explicitly cast as the violator of a shared moral value.
    \item \textbf{Identity-defense grammar} — the speaker invokes ingroup virtue vocabulary (\textit{tradition}, \textit{purity}, \textit{order}, \textit{loyalty}) to justify the attack, and simultaneously marks the target as a threat to those virtues (\textit{cult}, \textit{brainwashing}, \textit{foreign influence}), yielding a \emph{positive} bias on the MFT axes together with a \emph{negative} bias on the Realistic/Symbolic Threat axes. The target is cast as a danger to cherished values rather than as a moral violator.
\end{itemize}
A positive MFT score therefore does not exonerate a text; it indicates that the hostility is framed as defense of something cherished, rather than as the indictment of a wrongdoer. A central empirical claim of this chapter is that the three language communities systematically prefer different grammars — English skews toward direct condemnation, while Chinese and Japanese skew toward identity defense — so the apparently paradoxical mix of positive MFT scores and negative threat scores in the Chinese and Japanese rows of the figures below is the expected signature of identity-defense framing, not an inconsistency in the data.

%---------------------------------------------------------------
\section{Moral Framing by Language}
\label{sec:rq3_lang}

Figure~\ref{fig:rq3_lang_heatmap} (heatmap) and Figure~\ref{fig:rq3_lang_line} (line chart) display the mean bias score for each language on each moral axis. Together they tell a clear story: the three languages do not attack religion in the same way — they each seem to reach for a different moral vocabulary.

\begin{figure}[htbp]
    \centering
    \includegraphics[width=\linewidth]{RQ3/rq3_lang_heatmap.png}
    \caption{Mean moral bias scores per language across all seven axes. Warmer (red) = document language closer to the vice/threat pole of the axis; cooler (blue) = closer to the virtue/safety/cohesion pole. Note that in the identity-defense grammar (Chinese, Japanese; see reading guide above), positive values on the five classical MFT axes signal ingroup-virtue invocation and typically co-occur with negative values on the threat axes — both are markers of hostility.}
    \label{fig:rq3_lang_heatmap}
\end{figure}

\begin{figure}[htbp]
    \centering
    \includegraphics[width=\linewidth]{RQ3/rq3_lang_line.png}
    \caption{Per-axis mean bias scores for English, Chinese, and Japanese, shown as a parallel-coordinate profile. Each line traces the moral ``fingerprint'' of one language community. See the reading guide above for how to interpret the sign of each axis.}
    \label{fig:rq3_lang_line}
\end{figure}

\paragraph{English: a broad negative sweep across moral foundations.}
English discourse sits firmly in negative territory on the five classical Moral Foundation axes. The mean bias scores are $-0.081$ (Authority), $-0.074$ (Harm), $-0.065$ (Sanctity), $-0.063$ (Fairness), and $-0.050$ (Loyalty). In plain terms, English anti-religious texts tend to paint the target as harmful, unfair, corrupt, and authoritarian — a wide-ranging moral indictment. This pattern is consistent with what one might call \emph{moral condem\allowbreak nation-style attack}: the speaker argues that religion (or religious actors) is bad across many ethical dimensions at once.

Interestingly, English is the \emph{only} language where the Realistic Threat axis is positive ($+0.018$), meaning English texts are \emph{not} particularly prone to framing religion as a concrete physical danger.

\paragraph{Chinese and Japanese: threat framing stands out.}
Chinese and Japanese show a markedly different profile. On the five classical Moral Foundation axes (Harm, Fairness, Loyalty, Authority, Sanctity), both languages tend toward \emph{positive} scores — Chinese: $+0.177, +0.182, +0.108, +0.117, +0.168$; Japanese: $+0.069, +0.085, +0.024, \allowbreak +0.034, +0.077$. This means that, on average, Chinese- and Japanese-language texts do \emph{not} simply attack religion on classical moral grounds to the same degree as English.

However, both languages flip sharply negative on the two threat axes: Symbolic Threat ($-0.047$ for Chinese, $-0.045$ for Japanese) and Realistic Threat ($-0.012$ for Chinese, $-0.018$ for Japanese). This suggests that Chinese and Japanese anti-religious discourse leans into \emph{threat-based framing} — portraying religion as a danger to cultural identity, social cohesion, or physical safety — rather than mounting a broad moral condemnation.

A concise way to summarise this divergence: English discourse is more negatively aligned with harm-, fairness-, authority-, and sanctity-related frames, whereas Chinese and Japanese discourse shows its main negative alignment on the realistic-threat and symbolic-threat axes. The attack is present in all three languages, but the rhetorical weapon of choice differs substantially.


%---------------------------------------------------------------
\section{Moral Framing by Topic}
\label{sec:rq3_topic}

Beyond language, we ask whether certain \emph{topics} systematically attract more aggressive moral framing than others, and whether the \emph{type} of aggression (condemnation vs.\ threat) also varies across topics.

\subsection*{Heatmap: which topics engage which moral axes?}

Figure~\ref{fig:rq3_topic_heatmap} shows the mean bias score for each (topic, axis) combination across all languages.

\begin{figure}[htbp]
    \centering
    \includegraphics[width=\linewidth]{RQ3/rq3_topic_heatmap.png}
    \caption{Per-topic mean moral bias scores across all seven axes. Most topics show at least moderate negative alignment on Symbolic Threat, while several also engage Harm, Fairness, Authority, and Sanctity frames.}
    \label{fig:rq3_topic_heatmap}
\end{figure}

The dominant pattern in the heatmap is Symbolic Threat: the vast majority of topics show a negative bias on this axis, suggesting that framing religion or religious groups as a threat to cultural identity and social values is a near-universal feature of anti-religious discourse, regardless of the specific topic. Beyond that baseline, a subset of topics additionally engages the classical Moral Foundation axes — Harm, Fairness, Loyalty, Authority, and Sanctity — producing stronger and more diverse moral indictments.

\subsection*{Scatter plot: how hostile is each topic, and what kind of hostility?}

Figure~\ref{fig:rq3_topic_scatter} maps every topic as a bubble onto a two-dimensional moral space derived from Principal Component Analysis (PCA) of the seven bias scores. The bubble size reflects the number of documents in that topic, making it easy to see which topics matter most.

\begin{figure}[htbp]
    \centering
    \includegraphics[width=\linewidth]{RQ3/rq3_topic_scatter.png}
    \caption{PCA-based moral space of topics. PC1 (63.7\% of variance) captures overall negativity/attack orientation; PC2 (16.9\% of variance) separates threat-based framing (top) from condemnation/harm-based framing (bottom). Bubble size is proportional to document count.}
    \label{fig:rq3_topic_scatter}
\end{figure}

Reading the plot is straightforward with a simple compass:
\begin{itemize}
    \item \textbf{Left} = more attack-oriented overall; \textbf{Right} = more normative or discussion-oriented.
    \item \textbf{Top-left} = attack through threat narratives; \textbf{Bottom-left} = attack through moral condemnation.
    \item \textbf{Larger bubbles} deserve more analytical attention because they represent more data.
\end{itemize}

PC1 explains 63.7\% of the total variance, meaning the single most important thing to know about a topic is simply \emph{how aggressive} it is. PC2 adds another 16.9\%, capturing the distinction between \emph{what kind} of aggression is deployed.

\paragraph{Most hostile topics.}
The leftmost topics — those furthest into attack territory — are \textit{T23 Historical Critiques of Christian Violence} (far-left, slightly upward), \textit{T42 Faith and Heresy Debates}, \textit{T46 Religious Beliefs and Social Controversies}, and \textit{T39 Religious Controversies and Social Critiques} (far-left, well below centre). Topics like \textit{T17 Pagan-Christian Conflicts}, \textit{T11 2017 Pakistan Religious Violence}, \textit{T4 Religion, Sexuality, and Violence}, and \textit{T6 Church Abuse Victims} are also prominently on the left side. The pattern is clear: topics centred on historical violence, heresy, social conflict, and abuse systematically generate more negatively charged discourse.

\paragraph{Most normative topics.}
On the right side, topics such as \textit{T29 Christian Principles and Criticisms}, \textit{T43 Missionaries and Confucianism}, \textit{T33 Healthcare Rights and Institutions}, \textit{T34 Church-State Relations and Charities}, and \textit{T2 Women and Ordination} appear more moderate. These topics tend to involve theological debate, institutional policy, or principled disagreement — discussions where hostile moral framing is present but not dominant.

\paragraph{Large bubbles worth highlighting.}
Among the high-volume topics, \textit{T1 Christians and Trump's Politics} (large, bottom-left) is both very common and strongly condemnation-oriented. \textit{T4 Religion, Sexuality, and Violence} (large, left-upper) is common and more threat-oriented. \textit{T6 Church Abuse Victims} (large, bottom-left) is common and predominantly harm/condemnation-oriented. By contrast, \textit{T0 Separation of Church \& State / Japan} — the single largest bubble — sits near the centre-right, indicating that despite its high frequency, it tends toward relatively measured, institutionally framed discussion rather than moral attack.

To summarise, not all religious topics attract equal hostility. The strongest attack-oriented framing clusters around themes of historical violence, heresy, and social-sexual conflict, while institutional and theological topics remain comparatively more civil. Within the hostile cluster, some topics are driven by threat narratives (upper-left) and others by moral condemnation (lower-left).


%---------------------------------------------------------------
\section{Cross-Language Variation Within Topics}
\label{sec:rq3_topic_lang}

The previous two sections established that language and topic each shape moral framing. Here we bring them together and ask: \emph{when the same topic is discussed in different languages, does the moral framing shift — and if so, how?} To answer this, we move from aggregate language or topic profiles to individual \emph{(topic, language)} units and examine them jointly in Figure~\ref{fig:rq3_lang_bar} and Figure~\ref{fig:rq3_topic_lang_scatter}.

\begin{figure}[h]
    \centering
    \includegraphics[width=\linewidth]{RQ3/rq3_lang_bar.png}
    \caption{Per-language mean moral bias scores grouped by topic. Topics are sorted by overall attack intensity. Bars show the attack salience of each language within each topic.}
    \label{fig:rq3_lang_bar}
\end{figure}

\begin{figure}[h]
    \centering
    \includegraphics[width=\linewidth]{RQ3/rq3_topic_lang_scatter.png}
    \caption{PCA moral space of all topic-language combinations. Each point represents one (topic, language) unit: English (blue), Chinese (orange), Japanese (green). Arrows show the direction of increasing positive bias for each moral axis. Points far from the origin indicate extreme moral profiles.}
    \label{fig:rq3_topic_lang_scatter}
\end{figure}

\paragraph{How to read Figure~\ref{fig:rq3_topic_lang_scatter}.}
Unlike Figure~\ref{fig:rq3_topic_scatter} — where each bubble was an entire topic averaged across all languages — here each point is a \emph{language-within-topic} unit (e.g., English-T42, Chinese-T42, Japanese-T42 are three separate points). The axes carry the same meaning as before: leftward along PC1 means more attack-oriented overall; upward along PC2 means the attack leans more toward threat framing, downward means it leans more toward moral condemnation. The arrows indicate the positive poles of each moral axis, so a point situated \emph{opposite} an arrow is negatively biased on that axis.

\paragraph{Two questions to ask for each topic.}
For any topic appearing in multiple languages, we can ask two distinct questions. First, \emph{how much does attack intensity differ across languages?} — measured by how far apart the three coloured points lie along the PC1 axis. Second, \emph{does the type of attack differ?} — measured by vertical separation along PC2. Some topics show only one kind of divergence; the most informative cases show both.

\paragraph{Topics where attack intensity diverges across languages.}
The starkest example is \textit{T39 Religious Controversies and Social Critiques}, where the English (PC1~$= -2.27$) and Chinese (PC1~$= -2.32$) points sit deep in attack territory, driven by a sweep of negative scores across all moral axes, while the pattern holds differently in Japanese (only two documents, limiting inference). A similar intensity gap appears in \textit{T9 Divine-Human Nature / Jesus}: English-language discussions are strongly attack-oriented (PC1~$= -1.17$) and dominant on the Harm axis, whereas Chinese-language discussions of the same theological topic score positively on all five classical axes — the negative signal there comes mainly from Symbolic Threat (${-0.032}$). This illustrates that a topic's location in the heatmap of Section~\ref{sec:rq3_topic} is a cross-language average; the underlying language-specific profiles can diverge substantially.

\paragraph{Topics where attack type diverges, even when intensity is similar.}
\textit{T42 Faith and Heresy Debates} is the clearest example of within-topic divergence in the \emph{kind} of moral attack. Both Chinese (PC1~$= -2.38$) and Japanese (PC1~$= -2.80$) communities treat this topic as highly hostile — the intensity is comparable. But the dominant negative axis differs: Japanese discourse is most strongly negative on Symbolic Threat ($z = -1.82$), framing heresy primarily as a threat to cultural and religious identity, while Chinese discourse loads most heavily on Authority ($z = -1.56$), framing it as a subversion of legitimate institutional order. The vertical spread of the two points in Figure~\ref{fig:rq3_topic_lang_scatter} — Japanese above, Chinese below the horizontal — reflects precisely this difference in attack type.

\paragraph{Topics that are consistent across languages.}
Not all topics diverge. \textit{T6 Church Abuse Victims} and \textit{T1 Christians and Trump's Politics} show relatively similar PC1 positions across all three languages (T6: EN~$= -0.87$, ZH~$= -1.17$, JP~$= -1.21$; T1: EN~$= -0.70$, ZH~$= -0.43$, JP~$= -0.41$), and in each language the dominant negative axis is either Loyalty or Symbolic Threat. This cross-linguistic agreement suggests that some topics carry a near-universal moral charge — the exposure of institutional abuse or the entanglement of religion with political power provokes similar moral reactions regardless of the speaker's cultural background.

\paragraph{A structural asymmetry across all topics.}
One systematic pattern runs through every topic: English-language texts consistently score negative on all five classical Moral Foundation axes, while Chinese- and Japanese-language texts in the same topic score \emph{positive} on those axes and push negative only on the threat dimensions. This means that within any given topic, the English data point tends to drift downward (toward the condemnation quadrant), while the Chinese and Japanese points cluster near or above the horizontal midline. The moral vocabulary difference established at the language level in Section~\ref{sec:rq3_lang} is therefore not an artefact of topic composition — it persists even when we hold the topic constant.


%---------------------------------------------------------------
\section{Summary: Why Does Anti-Religious Hate Speech Happen?}
\label{sec:rq3_summary}
At its core, the answer to RQ3 is this: people attack Christianity — and religion more broadly — because they perceive it as violating, or as threatening, the moral values they hold most dear. Anti-religious hate speech is not simply the expression of dislike; it is a moral argument. The speaker is not saying ``I dislike this''; they are saying ``this is wrong, this is harmful, this is a threat to us.'' The specific moral charge, however, depends heavily on who is speaking and what they are talking about.

For English-speaking audiences, the dominant grievance is moral condemnation. Christianity is attacked for causing harm, enabling injustice, undermining legitimate authority, or degrading human dignity — the Harm, Fairness, Authority, and Sanctity axes all pull negative simultaneously. The implied complaint is that the religion, or its institutions, have betrayed a wide set of ethical standards that society takes for granted. Topics like church abuse scandals or religion in politics exemplify this: the speaker's anger is grounded in the belief that religion has actively done something wrong.

For Chinese- and Japanese-speaking audiences, the moral charge shifts. Here the primary complaint is not that religion has done something wrong in a universal ethical sense, but that it poses a threat — to cultural identity, social cohesion, or group safety. The Symbolic Threat and Realistic Threat axes turn negative while the classical Moral Foundation axes remain relatively positive. The implied complaint is closer to: ``this religion does not belong here, and its presence endangers us.'' Topics like faith and heresy debates concentrate the most extreme versions of this framing in both languages, though Chinese discourse leans more strongly into identity-threat narratives while Japanese discourse incorporates more direct moral criticism.

In both cases, the hate speech follows a recognisable logic: the speaker identifies a moral value they care about, constructs an argument that the target religion either violates it or threatens it, and uses that construction as the justification for attack. What changes across languages and topics is not only \emph{which value} is invoked but also the \emph{grammar} of the moral argument itself: English-language discourse tends to indict religion as the violator of a shared ethical code (direct-condemnation grammar), while Chinese- and Japanese-language discourse tends to position religion as a danger to cherished ingroup virtues (identity-defense grammar). Neither grammar is more hostile than the other; they are different rhetorical routes to the same end. And the choice of both value and grammar is far from random — it reflects the broader cultural and social concerns of each community, turning anti-religious discourse into a lens through which those communities articulate what they believe a good society should look like.


% self-defined macro: include bibliography even when compiling a single chapter
\subfilebibliography
\end{document}
